\chapter{Uveitis}

\section*{Definition}
Uveitis is inflammation of the uveal tract (iris, ciliary body, choroid) classified anatomically as anterior (iritis, iridocyclitis), intermediate (pars planitis), posterior (choroiditis, retinitis), or panuveitis.

\section*{What Happens in Your Body}
Breakdown of the blood-ocular barrier allows inflammatory cells and proteins to enter the anterior chamber, vitreous, or retina. Immune signaling proteins-mediated infiltration (T-cell driven in most cases) and release of prostaglandins and other inflammatory mediators produce the clinical features of pain, photophobia, and cellular reaction.

\section*{Causes}
\begin{itemize}
  \item \textbf{Autoimmune:} Ankylosing spondylitis, reactive arthritis, psoriatic arthritis, juvenile idiopathic arthritis, Behcet disease, sarcoidosis
  \item \textbf{Infectious:} Herpes simplex, herpes zoster, cytomegalovirus, toxoplasmosis, tuberculosis, syphilis, Lyme disease
  \item \textbf{Masquerade syndromes:} Intraocular lymphoma, retinoblastoma, uveal melanoma
  \item \textbf{idiopathic:} No identifiable cause in up to 50\% of anterior uveitis cases
  \item \textbf{Traumatic:} Post-surgical or blunt trauma (sympathetic ophthalmia)
\end{itemize}

\section*{When to See a Doctor}
See your doctor if:
\begin{itemize}
  \item Unilateral eye pain, photophobia, and redness with blurred vision
  \item Small or irregular pupil (posterior synechiae formation)
  \item Floaters or decreased vision without pain (suggests intermediate or posterior uveitis)
  \item Recurrent episodes of uveitis in the same or fellow eye
  \item Known systemic autoimmune disease with new ocular symptoms
\end{itemize}

\section*{Related Symptoms}
\begin{itemize}
  \item Ocular redness (ciliary flush)
  \item Photophobia
  \item Blurred vision and floaters
  \item Eye pain (deep, aching, worse with accommodation)
  \item Tearing and blepharospasm
\end{itemize}
\textbf{Important:} Anterior uveitis is the most common form (90\% of cases) and is frequently associated with HLA-B27 positivity; prompt treatment with topical steroids prevents synechiae and glaucoma.

\section*{Self-Care / Home Management}
\begin{itemize}
  \item Sunglasses to reduce photophobia
  \item Warm compresses for comfort
  \item Avoidance of contact lens wear during active inflammation
  \item Compliance with prescribed steroid and cycloplegic eye drops
  \item Monitoring for symptoms of increased IOP (pain, halos, worsening vision) while on topical steroids
\end{itemize}

\section*{How Your Doctor Will Evaluate You}
\begin{itemize}
  \item Slit-lamp examination for anterior chamber cells and flare (SUN grading system)
  \item Keratic precipitates (granulomatous or non-granulomatous)
  \item Fundus examination with dilated pupil for vitreous cells, snowballs, retinal infiltrates, or vasculitis
  \item IOP measurement (may be elevated or low)
  \item Fluorescein angiography or OCT for posterior uveitis
  \item Laboratory workup: CBC, ESR, CRP, HLA-B27, syphilis serology, quantiferon-TB, ACE level, and lysozyme (individualized)
\end{itemize}

\section*{Treatment Options}
\begin{itemize}
  \item \textbf{Anterior:} Topical prednisolone acetate 1\% (frequent dosing tapered slowly) and cycloplegic (homatropine or atropine)
  \item \textbf{Intermediate/posterior:} Periocular corticosteroid injection (triamcinolone) or systemic steroids
  \item \textbf{Immunomodulatory therapy:} Methotrexate, mycophenolate mofetil, azathioprine, or biologic agents (adalimumab, infliximab)
  \item \textbf{Infectious:} Directed antimicrobial therapy (e.g., valganciclovir for CMV, trimethoprim-sulfamethoxazole for toxoplasmosis)
  \item \textbf{Complications:} Management of secondary glaucoma, cataract, cystoid macular swelling, and band keratopathy
  \item \textbf{Long-term monitoring:} Serial examinations to detect progression, complications, or disease recurrence
\end{itemize}

\section*{References}
\begin{thebibliography}{99}
\bibitem{uveitis_symptoms:ref1} Jabs DA, Nussenblatt RB, Rosenbaum JT, et al. ``Standardization of uveitis nomenclature (SUN) for reporting clinical data. Results of a first international workshop.'' \textit{American Journal of Ophthalmology}, 2005;140(3):509--516.
\bibitem{uveitis_symptoms:ref2} Guly CM, Forrester JV. ``Uveitis.'' \textit{BMJ}, 2010;341:c4376.
\bibitem{uveitis_symptoms:ref3} American Academy of Ophthalmology. ``Intraocular Inflammation and Uveitis.'' \textit{Basic and Clinical Science Course}. AAO; 2023.
\bibitem{uveitis_symptoms:ref4} Thorne JE, Suhler E, et al. ``Uveitis: clinical features and diagnosis.'' \textit{UpToDate}, Post TW (Ed), UpToDate, Waltham, MA; 2025.
\bibitem{uveitis_symptoms:ref5} Acharya NR, Tham VM, Johnston WH, et al. ``Etiologies of uveitis in a referral center.'' \textit{Eye}, 2024;38(2):294--300.
\bibitem{uveitis_symptoms:ref6} Jameson JL, et al. \textit{Harrison\'s Principles of Internal Medicine}. 21st ed. McGraw-Hill; 2022. Chapter 29: Ocular Disorders.
\end{thebibliography}
